This study answers the two research questions posed in Section~\ref{sec:intro}, showing that publication-based topic representations can support researcher-to-research-group recommendation. BERTopic produced the numerically best ranking performance on every metric, though its margin over LDA was statistically significant on only 3 of 13 tested metrics (Section~\ref{sec:results}), while NMF was most efficient. The main contribution is a BERTopic-based framework representing publications, researchers, and research groups in a shared topic space; its novelty lies in the controlled comparison of LDA, NMF, and BERTopic rather than in a new algorithm, clarifying how topic representation affects ranking quality and showing that a numerical ranking alone does not establish which differences are statistically dependable.

Topic coherence is not a reliable substitute for downstream evaluation: NMF generated the most coherent topics but the weakest recommendations, while BERTopic's better ranking quality came at a much higher cost and may partly reflect its externally pretrained SPECTER embeddings rather than an intrinsic property of contextual topic modeling. Model selection should reflect intended use: BERTopic when quality has priority and offline training is acceptable, LDA as a reasonable default at a fraction of the cost when that fragile margin does not justify the expense, and NMF when frequent retraining or limited resources dominate, at the cost of recommendation quality.

The evaluation is limited to CSRankings data from 12 computer science research groups, one researcher-level split, and profiles derived only from titles and abstracts; it also does not separate BERTopic's contextual architecture from SPECTER's external pretraining, so results should be read as topic-modeling pipelines as commonly deployed rather than an ablation of the contextual mechanism. Uneven group sizes and publication counts may further affect profile quality, particularly for the smallest groups, limiting generalization beyond this dataset. The significance tests in Section~\ref{sec:results} quantify uncertainty within this split; future work should evaluate repeated splits, extend to additional institutions or disciplines, examine performance by profile size, test citation, collaboration, and temporal information and alternative aggregation, and ablate BERTopic's components to determine whether its advantage over LDA persists once the pretraining confound is controlled.
